\begin{helper}
Let $n\ge0$ be an integer, let $k\in\mathbb{Z}$ and let $c',d'$ be real with
$c'>0$.  Write
\[
  \mathcal{A}=\{\zeta\in\mathbb{C}:c'\le|\zeta|\le d'\}
\]
for the compact annulus between the two radii, and for
$F\colon\mathbb{C}\to\mathbb{C}$ write
$\widehat F_k(s)=\frac1{2\pi}\int_0^{2\pi}F(se^{i\theta})e^{i(n-k)\theta}d\theta$
for its $k$-th Fourier coefficient at radius $s$, as in
\noderef{FourierCoefOfH}.  Then the operator $F\mapsto\widehat F_k$ has the
following two properties.
\begin{enumerate}
\item[(a)] If $W\colon\mathbb{C}\to\mathbb{C}$ is measurable and integrable on
  $\mathcal{A}$, then $s\mapsto\widehat W_k(s)$ is integrable on $(c',d')$,
  and for almost every $s\in(c',d')$ the circle restriction
  $\theta\mapsto W(se^{i\theta})$ is interval-integrable over $[0,2\pi]$.
\item[(b)] The coefficient depends continuously on the function in $L^1$: if
  $F_l,F\colon\mathbb{C}\to\mathbb{C}$ are measurable and integrable on
  $\mathbb{C}$ and $\int_{\mathbb{C}}|F_l-F|\,d\sigma\to0$, then
  \[
    \int_{(c',d')}\bigl\|\widehat{(F_l)}_k(s)-\widehat F_k(s)\bigr\|\,ds
    \longrightarrow0,
  \]
  the integral being the $[0,\infty]$-valued Lebesgue integral.
\end{enumerate}

Both are properties of the one operator $F\mapsto\widehat F_k$, which is why
they are stated together: (b) is proved from the same estimate that proves
(a), and (a) is what makes the inner integral defining $\widehat F_k$
meaningful in the first place.

This is the measure-theoretic foundation under the circle-average argument.
The paper reads $a_k\in C[0,1)\cap C^\infty(0,1)$ off its standing smoothness
hypothesis at \texttt{paper.tex:265-266}; the tablet carries no such
hypothesis and replaces that step by \noderef{CircleAvgLocAbsCont}, which
needs exactly these two facts before it can begin --- (a) to know that
$\widehat W_k$ is a genuine integrable function of the radius, and (b) to pass
to the limit in a mollification, where the approximating functions converge to
$W$ only in $L^1$.

The hypothesis $c'>0$ is essential and is used twice: it keeps the Jacobian
$s$ of the polar map bounded below on $\mathcal{A}$, which is what converts an
$L^1$ bound on the plane into an $L^1$ bound in the polar variables, and it
bounds the factor $1/s$ implicit in the pointwise estimate behind (b).  The
upper radius $d'$ plays no such role; it only makes the annulus bounded.
\end{helper}
\begin{proof}
Throughout write
\[
  \mathrm{ci}(s,\theta)=se^{i\theta},\qquad
  \mathrm{wt}(\theta)=e^{i(n-k)\theta},
\]
and $d\sigma$ for plane Lebesgue measure.  By \noderef{FourierCoefOfH},
$\widehat F_k(s)=\frac1{2\pi}\int_0^{2\pi}F(\mathrm{ci}(s,\theta))\,
\mathrm{wt}(\theta)\,d\theta$.  Note $|\mathrm{wt}(\theta)|=1$ and
$\mathrm{wt}(\theta)\neq0$ for every $\theta$, since $n-k$ is an integer and
$\theta$ is real, so the exponent is purely imaginary; and
$|\mathrm{ci}(s,\theta)|=|s|$.  Both $\mathrm{ci}$ and $\mathrm{wt}$ are
continuous.  The set $\mathcal{A}$ is closed and bounded, hence compact, hence
measurable.

\emph{Step 0: two standard transfers.}  First, the polar change of variables:
for every measurable $G\colon\mathbb{C}\to[0,\infty]$,
\begin{equation}\tag{1}
  \int_{(0,\infty)\times(-\pi,\pi)}s\,G\bigl(\mathrm{ci}(s,\theta)\bigr)
    \,d(s,\theta)
  =\int_{\mathbb{C}}G\,d\sigma ,
\end{equation}
the map $(s,\theta)\mapsto se^{i\theta}$ being a diffeomorphism of
$(0,\infty)\times(-\pi,\pi)$ onto the plane minus a null ray, with Jacobian
$s$.  Second, the window change: for fixed $s$ the function
$\theta\mapsto F(\mathrm{ci}(s,\theta))\mathrm{wt}(\theta)$ is $2\pi$-periodic,
because $\mathrm{ci}$ and $\mathrm{wt}$ are, so
\begin{equation}\tag{2}
  \int_0^{2\pi}F\bigl(\mathrm{ci}(s,\theta)\bigr)\mathrm{wt}(\theta)\,d\theta
  =\int_{(-\pi,\pi)}F\bigl(\mathrm{ci}(s,\theta)\bigr)\mathrm{wt}(\theta)
    \,d\theta ,
\end{equation}
a period's worth of integral being independent of where the period starts.

\emph{Step 1: the key lower bound, where $c'>0$ enters.}  For every
$F\colon\mathbb{C}\to\mathbb{C}$,
\begin{equation}\tag{3}
  c'\int_{(c',d')\times(-\pi,\pi)}
      \bigl\|F\bigl(\mathrm{ci}(s,\theta)\bigr)\bigr\|\,d(s,\theta)
  \le\int_{\mathbb{C}}\|F\|\,d\sigma ,
\end{equation}
all integrals being $[0,\infty]$-valued.  Indeed on $(c',d')\times(-\pi,\pi)$
we have $s>c'$, so pointwise
$c'\|F(\mathrm{ci}(s,\theta))\|\le s\,\|F(\mathrm{ci}(s,\theta))\|$; hence
the left side is at most
$\int_{(c',d')\times(-\pi,\pi)}s\|F(\mathrm{ci}(s,\theta))\|$, which is at most
the same integral over the larger set $(0,\infty)\times(-\pi,\pi)$ because
$(c',d')\subseteq(0,\infty)$ and the integrand is nonnegative, and that equals
$\int_{\mathbb{C}}\|F\|$ by (1) applied to $G=\|F\|$.

\emph{Step 2: the product is integrable on the rectangle.}  Let $F$ be
measurable and integrable on $\mathbb{C}$.  Then
$(s,\theta)\mapsto F(\mathrm{ci}(s,\theta))\mathrm{wt}(\theta)$ is measurable
on $(c',d')\times(-\pi,\pi)$, being the product of the composition of the
measurable $F$ with the continuous $\mathrm{ci}$ and of the continuous
$\mathrm{wt}$; and it is integrable there, because
$\|F(\mathrm{ci}(s,\theta))\mathrm{wt}(\theta)\|
=\|F(\mathrm{ci}(s,\theta))\|$ by $|\mathrm{wt}|=1$, and by (3)
\begin{equation}\tag{4}
  \int_{(c',d')\times(-\pi,\pi)}
    \bigl\|F\bigl(\mathrm{ci}(s,\theta)\bigr)\bigr\|\,d(s,\theta)
  \le\frac1{c'}\int_{\mathbb{C}}\|F\|\,d\sigma<\infty ,
\end{equation}
where dividing by $c'$ is legitimate because $c'>0$.

\emph{Step 3: proof of (a).}  Let $W$ be measurable and integrable on
$\mathcal{A}$, and put $F=\mathbf1_{\mathcal{A}}W$.  Then $F$ is measurable,
$\mathcal{A}$ being measurable, and $F$ is integrable on all of $\mathbb{C}$,
since it vanishes off $\mathcal{A}$ and agrees with $W$ there.  So Step 2
applies to $F$, and Fubini's theorem for integrable functions on the product
$(c',d')\times(-\pi,\pi)$ gives both
\begin{equation}\tag{5}
  \text{for a.e. }s\in(c',d'):\quad
  \theta\mapsto F\bigl(\mathrm{ci}(s,\theta)\bigr)\mathrm{wt}(\theta)
  \text{ is integrable on }(-\pi,\pi),
\end{equation}
and
\begin{equation}\tag{6}
  s\mapsto\int_{(-\pi,\pi)}F\bigl(\mathrm{ci}(s,\theta)\bigr)
    \mathrm{wt}(\theta)\,d\theta
  \text{ is integrable on }(c',d').
\end{equation}
Now let $s\in(c',d')$.  For every $\theta$ we have
$|\mathrm{ci}(s,\theta)|=|s|=s\in(c',d')\subseteq[c',d']$, so
$\mathrm{ci}(s,\theta)\in\mathcal{A}$ and therefore
$F(\mathrm{ci}(s,\theta))=W(\mathrm{ci}(s,\theta))$.  Hence, by the definition
of the coefficient and by (2),
\[
  \widehat W_k(s)
  =\frac1{2\pi}\int_0^{2\pi}W\bigl(\mathrm{ci}(s,\theta)\bigr)
     \mathrm{wt}(\theta)\,d\theta
  =\frac1{2\pi}\int_{(-\pi,\pi)}F\bigl(\mathrm{ci}(s,\theta)\bigr)
     \mathrm{wt}(\theta)\,d\theta
  \qquad\bigl(s\in(c',d')\bigr),
\]
so on $(c',d')$ the function $s\mapsto\widehat W_k(s)$ is the constant
$\frac1{2\pi}$ times the function of (6), and is therefore integrable on
$(c',d')$.  That is the first half of (a).

For the second half, fix an $s\in(c',d')$ for which (5) holds.  Dividing by
$\mathrm{wt}$, which is continuous and never zero, the function
$\theta\mapsto F(\mathrm{ci}(s,\theta))$ is integrable on $(-\pi,\pi)$; it
equals $\theta\mapsto W(\mathrm{ci}(s,\theta))$, as just observed; and that
function is $2\pi$-periodic, so by the same period-shift argument as in (2) it
is interval-integrable over $[0,2\pi]$.  Since the exceptional set of $s$ is
null, this holds for almost every $s\in(c',d')$, which is the second half of
(a).

\emph{Step 4: the pointwise estimate behind (b).}  Let $G_1,G_2$ be measurable
and integrable on $\mathbb{C}$ and let $s$ be real.  The coefficient is linear
in the function, so
$\widehat{(G_1)}_k(s)-\widehat{(G_2)}_k(s)
=\frac1{2\pi}\int_0^{2\pi}(G_1-G_2)(\mathrm{ci}(s,\theta))
\mathrm{wt}(\theta)\,d\theta$, and hence, by the triangle inequality for the
integral, by $|\mathrm{wt}|=1$, and by (2),
\begin{equation}\tag{7}
  \bigl\|\widehat{(G_1)}_k(s)-\widehat{(G_2)}_k(s)\bigr\|
  \le\frac1{2\pi}\int_{(-\pi,\pi)}
      \bigl\|(G_1-G_2)\bigl(\mathrm{ci}(s,\theta)\bigr)\bigr\|\,d\theta .
\end{equation}

\emph{Step 5: proof of (b).}  Let $F_l$, $F$ be as in (b) and put
$\eta_l=\int_{\mathbb{C}}\|F_l-F\|\,d\sigma$, so $\eta_l\to0$ by hypothesis.
Integrating (7) with $G_1=F_l$, $G_2=F$ over $s\in(c',d')$, using Tonelli's
theorem to read the iterated integral of a nonnegative measurable function as
an integral over the rectangle, and then applying (3) to $F_l-F$,
\[
  \int_{(c',d')}\bigl\|\widehat{(F_l)}_k(s)-\widehat F_k(s)\bigr\|\,ds
  \le\frac1{2\pi}\int_{(c',d')\times(-\pi,\pi)}
      \bigl\|(F_l-F)\bigl(\mathrm{ci}(s,\theta)\bigr)\bigr\|\,d(s,\theta)
  \le\frac1{2\pi}\cdot\frac{\eta_l}{c'} .
\]
Both constants $\frac1{2\pi}$ and $\frac1{c'}$ are finite and independent of
$l$, the first because $\pi>0$ and the second because $c'>0$, so the
right-hand side tends to $0$ as $l\to\infty$.  The left-hand side is
nonnegative, so it is squeezed between $0$ and a sequence tending to $0$ and
therefore tends to $0$ itself.  That is (b), and the proof is complete.
\end{proof}
